\chapter*{Заключение}
\addcontentsline{toc}{chapter}{Заключение}

В рамках работы была исследована задача построения вопросно-ответного ассистента для русскоязычной корпоративной базы знаний RobyMAX. В аналитической части были рассмотрены особенности предметной области, проблема поиска по разнородным документам и ограничения ручного и полнотекстового поиска. На основе анализа был обоснован выбор базового Retrieval-Augmented Generation как подхода, позволяющего совмещать семантический поиск, генерацию ответа и возврат источников.

В теоретической части были рассмотрены большие языковые модели, embedding-представления текста, векторный поиск и общая схема RAG. Было показано, что использование внешней базы знаний позволяет не дообучать LLM под каждый новый набор документов, а передавать модели релевантный контекст, найденный в актуальном корпусе.

В инженерной части была спроектирована архитектура базового RAG-ассистента. Она включает offline-этап подготовки корпуса и online-этап ответа на вопрос. Offline-этап состоит из регистрации документов, парсинга, нормализации, чанкинга, построения embeddings и индексации в Qdrant. Online-этап включает кодирование вопроса, dense retrieval, сбор контекста, генерацию ответа и возврат источников.

В практической части был описан реализованный MVP-пайплайн. Корпус включает 36 документов форматов DOCX, PDF и XLSX; после парсинга получено 2300 структурированных блоков и 481 чанк. Для построения embeddings используется модель \texttt{deepvk/USER-bge-m3}, для хранения и поиска векторов — Qdrant, для генерации ответа — локальная модель \texttt{qwen2.5:7b-instruct}. Первичная retrieval-проверка базового dense search на 150 вопросах показала \texttt{Hit@5=0.7467}, \texttt{MRR@5=0.5676} и \texttt{Recall@5=0.6522}. Эти результаты подтверждают работоспособность минимального сценария, но не означают завершенность системы как production-ready решения.

Таким образом, поставленная цель достигнута: был разработан и описан базовый прототип RAG-ассистента для поиска информации в корпусе документов RobyMAX и генерации ответа по найденному контексту. Полученная MVP-версия является основой для дальнейшего развития, поскольку реализует полный путь от документов до ответа и позволяет отдельно оценивать качество подготовки данных, поиска и генерации.

Дальнейшее развитие целесообразно вести по следующим направлениям:
\begin{compactenum}
  \item добавление BM25 и гибридного retrieval;
  \item добавление semantic reranker;
  \item подключение LangGraph workflow;
  \item добавление domain gate и проверки достаточности контекста;
  \item расширение RAG evaluation;
  \item создание API и пользовательского интерфейса;
  \item оптимизация производительности и контейнеризация.
\end{compactenum}

%%% Local Variables:
%%% TeX-engine: xetex
%%% End:
